\section{The mathematical definition}\label{sec:10-modules:01-definition:1}

Given a ring \(R\) (Chapter 8), a (left) \textbf{\(R\)-module} is an abelian group \((M, +, 0, -(-))\) together with a scalar action \(R \times M \to M\), written \(r \cdot m\), satisfying:

\[
\begin{aligned}
\text{(M1)}&\quad r \cdot (m + n) = r\cdot m + r \cdot n \\
\text{(M2)}&\quad (r + s) \cdot m = r \cdot m + s \cdot m \\
\text{(M3)}&\quad (r \cdot s) \cdot m = r \cdot (s \cdot m) \\
\text{(M4)}&\quad 1 \cdot m = m
\end{aligned}
\]

for all \(r, s \in R\), \(m, n \in M\). This is exactly the vector space definition, with ``field'' replaced by ``ring''. That extra generality is the whole point: \(\mathbb{Z}\)-modules are abelian groups, and \(k[x]\)-modules (for \(k\) a field) are vector spaces equipped with a chosen linear endomorphism. Most importantly for this book, a representation of a quiver \(Q\) is precisely a module over the path algebra \(kQ\); this is why the present chapter is placed immediately before Chapter 11.
